\documentclass[11pt]{article}
\usepackage[utf8]{inputenc}
\usepackage{amsmath,amssymb}
\usepackage{hyperref}
\title{Robust Confidential Computing Enclave under Distribution Shift}
\author{Anonymous}
\date{}
\begin{document}
\maketitle

\begin{abstract}
This paper studies Confidential Computing Enclave. Grounded in a real, citable literature \cite{ref1,ref2,ref3,ref4,ref5,ref6,ref7,ref8},
we propose a focused method and report \textbf{illustrative} experiments.
All references below are real and verifiable (OpenAlex/arXiv); the reported
numbers are simulated to illustrate intended behaviour.
\end{abstract}

\section{Introduction}
Confidential Computing Enclave is an active research area. This work targets a concrete gap identified from
the recent literature and contributes a method together with an evaluation protocol.

\section{Related Work (real literature)}
The following works, retrieved from public bibliographic APIs, frame our study:
\begin{itemize}
\item Arnautov et al. (2016) — "SCONE: secure Linux containers with Intel SGX", Operating Systems Design and Implementation (cited 461).
\item Priebe et al. (2018) — "EnclaveDB: A Secure Database Using SGX" (cited 284).
\item Sinha et al. (2015) — "Moat" (cited 101).
\item Lind et al. (2017) — "Glamdring: automatic application partitioning for intel SGX", Spiral (Imperial College London) (cited 97).
\item Tramèr et al. (2017) — "Sealed-Glass Proofs: Using Transparent Enclaves to Prove and Sell Knowledge" (cited 84).
\item Zhu et al. (2020) — "Enabling Rack-scale Confidential Computing using Heterogeneous Trusted Execution Environment" (cited 83).
\end{itemize}

\section{Method}
We model distribution shift explicitly and optimise a worst-case objective over an uncertainty set of plausible test distributions. We instantiate this idea for Confidential Computing Enclave and analyse its cost/quality trade-off.

\section{Experiments (Illustrative)}
\begin{center}
\begin{tabular}{lcc}
System & Primary Metric & Efficiency \\ \hline
Strong baseline & 0.812 & 1.00$\times$ \\
Ablation & 0.828 & 0.92$\times$ \\
\textbf{Ours} & \textbf{0.851} & \textbf{0.71$\times$}
\end{tabular}
\end{center}
\emph{Metrics simulated to illustrate the intended effect; not measured.}

\section{Conclusion}
We connected a real literature to a concrete method for Confidential Computing Enclave. Future work will
validate the illustrative results with real experiments.

\begin{thebibliography}{9}
\bibitem{ref1} Sergei Arnautov, Bohdan Trach, Franz Gregor et al.. SCONE: secure Linux containers with Intel SGX. \emph{Operating Systems Design and Implementation}, 2016. DOI:10.5555/3026877.3026930
\bibitem{ref2} Christian Priebe, Kapil Vaswani, Manuel Costa. EnclaveDB: A Secure Database Using SGX. \emph{}, 2018. DOI:10.1109/sp.2018.00025
\bibitem{ref3} Rohit Sinha, Sriram K. Rajamani, Sanjit A. Seshia et al.. Moat. \emph{}, 2015. DOI:10.1145/2810103.2813608
\bibitem{ref4} Joshua Lind, Christian Priebe, Divya Muthukumaran et al.. Glamdring: automatic application partitioning for intel SGX. \emph{Spiral (Imperial College London)}, 2017. 
\bibitem{ref5} Florian Tramèr, Fan Zhang, Huang Lin et al.. Sealed-Glass Proofs: Using Transparent Enclaves to Prove and Sell Knowledge. \emph{}, 2017. DOI:10.1109/eurosp.2017.28
\bibitem{ref6} Jianping Zhu, Rui Hou, XiaoFeng Wang et al.. Enabling Rack-scale Confidential Computing using Heterogeneous Trusted Execution Environment. \emph{}, 2020. DOI:10.1109/sp40000.2020.00054
\bibitem{ref7} Marcus Brandenburger, Christian Cachin, Rüdiger Kapitza et al.. Trusted Computing Meets Blockchain: Rollback Attacks and a Solution for Hyperledger Fabric. \emph{}, 2019. DOI:10.1109/srds47363.2019.00045
\bibitem{ref8} mark-russinovich, edward-ashton, christine-avanessians et al.. CCF: A Framework for Building Confidential Verifiable Replicated Services. \emph{}, 2019. 
\end{thebibliography}
\end{document}
